\documentclass[11pt,a4paper]{article}
\usepackage{amsmath,amssymb,amsthm,mathtools}
\usepackage[margin=1in]{geometry}
\usepackage{hyperref}

%% Theorem environments
\newtheorem{theorem}{Theorem}[section]
\newtheorem{lemma}[theorem]{Lemma}
\newtheorem{proposition}[theorem]{Proposition}
\newtheorem{corollary}[theorem]{Corollary}
\newtheorem{claim}[theorem]{Claim}
\theoremstyle{definition}
\newtheorem{definition}[theorem]{Definition}
\newtheorem{remark}[theorem]{Remark}

%% Notation shortcuts
\newcommand{\T}{\mathbb{T}}
\newcommand{\R}{\mathbb{R}}
\newcommand{\E}{\mathbb{E}}
\newcommand{\Dc}{\mathcal{D}}
\newcommand{\Hc}{\mathcal{H}}
\newcommand{\norm}[1]{\left\|#1\right\|}
\newcommand{\wick}[1]{:\!#1\!:}
\newcommand{\ip}[2]{\langle #1,\, #2 \rangle}

\title{\textbf{Problem 1: Equivalence of $\Phi^4_3$ Measure under Shift}\\[4pt]
\large Quasi-invariance via Cameron--Martin theory and Wick polynomial expansion}
\author{}
\date{}

\begin{document}
\maketitle

%% ================================================================
\section*{Answer and Construction}
%% ================================================================

\noindent\textbf{Answer: Yes.}  The measures $\mu$ and $T_\psi^*\mu := (T_\psi)_*\mu$ are \emph{equivalent} (mutually absolutely continuous) for every smooth function $\psi:\T^3\to\R$.

Specifically, writing $\mu = Z^{-1}e^{-V}\,\mu_0$ for the Gaussian free field reference measure $\mu_0$ and the renormalized quartic potential $V$, the Radon--Nikodym derivative is
\begin{equation}\label{eq:RN_main}
\frac{d\,T_\psi^*\mu}{d\mu}(\phi)
\;=\;
\exp\!\Bigl(
  4\lambda\!\int_{\T^3}\!\psi\,\wick{\phi^3}\,dx
  -6\lambda\!\int_{\T^3}\!\psi^2\,\wick{\phi^2}\,dx
  +4\lambda\!\int_{\T^3}\!\psi^3\phi\,dx
  -\lambda\norm{\psi}_{L^4}^4
  +\phi\!\bigl((-\Delta+m^2)\psi\bigr)
  -\tfrac{1}{2}\norm{\psi}_{\Hc}^2
\Bigr),
\end{equation}
where $\wick{\phi^n}$ denotes the $n$-th Wick power with respect to $\mu_0$ and $\norm{\cdot}_\Hc^2 = \ip{\cdot}{(-\Delta+m^2)(\cdot)}_{L^2}$.
This derivative is strictly positive $\mu$-almost surely and satisfies $\E_\mu\bigl[\frac{d\,T_\psi^*\mu}{d\mu}\bigr]=1$.
The reverse absolute continuity $\mu\ll T_\psi^*\mu$ follows by the same argument applied to $-\psi$.

%% ================================================================
\section{Setup, Definitions, and Key Tools}
%% ================================================================

\subsection{Objects of the problem}

Let $\T^3 = (\R/\Z)^3$ be the three-dimensional unit torus, and let $\Dc'(\T^3)$ denote its space of distributions, equipped with the $\sigma$-algebra generated by evaluation maps $\phi\mapsto\phi(f)$ for $f\in C^\infty(\T^3)$.

\begin{definition}[Gaussian free field $\mu_0$]
Fix $m>0$ and $\lambda>0$.  The \emph{Gaussian free field} (GFF) $\mu_0$ on $\Dc'(\T^3)$ is the unique centred Gaussian measure with covariance
\[
\E_{\mu_0}[\phi(f)\,\phi(g)] = \ip{f}{(-\Delta+m^2)^{-1}g}_{L^2(\T^3)},
\qquad f,g\in C^\infty(\T^3).
\]
The operator $C := (-\Delta+m^2)^{-1}$ maps $L^2\to H^2$ by elliptic regularity.
The field $\phi$ is $\mu_0$-almost surely in $H^{-1/2-\varepsilon}(\T^3)$ for every $\varepsilon>0$, but not in $L^2$.
\end{definition}

\begin{definition}[Wick powers]\label{def:wick}
Let $\phi_\varepsilon = \rho_\varepsilon*\phi$ be the UV regularisation of $\phi$ with mollifier $\rho_\varepsilon$, and set $c_\varepsilon := \E_{\mu_0}[\phi_\varepsilon(x)^2]$ (which diverges as $\varepsilon\to 0$ in three dimensions).  The Wick-renormalised powers are
\begin{align*}
\wick{\phi_\varepsilon^2}(x) &:= \phi_\varepsilon(x)^2 - c_\varepsilon,\\
\wick{\phi_\varepsilon^3}(x) &:= \phi_\varepsilon(x)^3 - 3c_\varepsilon\,\phi_\varepsilon(x),\\
\wick{\phi_\varepsilon^4}(x) &:= \phi_\varepsilon(x)^4 - 6c_\varepsilon\,\phi_\varepsilon(x)^2 + 3c_\varepsilon^2.
\end{align*}
The integrated functionals $\int_{\T^3}\wick{\phi_\varepsilon^n}\,dx$ converge in $L^2(\mu_0)$ as $\varepsilon\to 0$ for $n\in\{2,3,4\}$, defining the renormalised Wick monomials $\int\wick{\phi^n}_{\mathrm{ren}}\,dx$; we write $\wick{\phi^n}$ for these limiting objects and omit the subscript ``ren'' henceforth.
\end{definition}

\begin{definition}[$\Phi^4_3$ measure $\mu$]\label{def:phi4}
The \emph{$\Phi^4_3$ measure} $\mu$ is the weak limit
\[
\mu := \lim_{\varepsilon\to 0}\mu_\varepsilon,
\]
where
\[
d\mu_\varepsilon(\phi) := Z_\varepsilon^{-1}\exp\!\Bigl(-\lambda\!\int_{\T^3}\wick{\phi_\varepsilon^4}\,dx
   -a_\varepsilon\!\int_{\T^3}\wick{\phi_\varepsilon^2}\,dx\Bigr)\,d\mu_0(\phi),
\]
with $a_\varepsilon$ a mass-renormalisation counterterm satisfying $a_\varepsilon\to+\infty$ at a rate determined by the UV divergences, and $Z_\varepsilon>0$ the normalising constant.
The limit $\mu$ exists as a probability measure on $\Dc'(\T^3)$~\cite{GlimmJaffe1987, Hairer2014, GIP2015}.
We write formally $d\mu = Z^{-1}e^{-V(\phi)}\,d\mu_0(\phi)$, where
\[
V(\phi) := \lambda\!\int_{\T^3}\wick{\phi^4}(x)\,dx
\]
is the renormalized quartic potential and $Z = \E_{\mu_0}[e^{-V}]\in(0,\infty)$.
\end{definition}

\subsection{Cameron--Martin theorem for $\mu_0$}

\begin{theorem}[Cameron--Martin]\label{thm:CM}
The Cameron--Martin space of $\mu_0$ is
\[
\Hc := C^{1/2}(L^2(\T^3)) = H^1(\T^3),
\]
equipped with the inner product $\ip{f}{g}_\Hc := \ip{f}{(-\Delta+m^2)g}_{L^2}$ and norm $\norm{f}_\Hc^2 = \norm{\nabla f}_{L^2}^2 + m^2\norm{f}_{L^2}^2$.

For any $\psi\in\Hc$, the measures $(T_\psi)_*\mu_0$ and $\mu_0$ are mutually absolutely continuous, and
\begin{equation}\label{eq:CM_RN}
R_\psi(\phi) :=\frac{d\,(T_\psi)_*\mu_0}{d\mu_0}(\phi)
= \exp\!\Bigl(\phi\bigl((-\Delta+m^2)\psi\bigr) - \tfrac{1}{2}\norm{\psi}_\Hc^2\Bigr),
\end{equation}
where $\phi(g)$ denotes the distributional pairing of $\phi\in\Dc'(\T^3)$ with $g\in C^\infty(\T^3)$.

\begin{proof}
This is the classical Cameron--Martin theorem~\cite[Theorem~2.4.3]{Bogachev1998} applied to the Gaussian measure $\mu_0$.
We verify the hypotheses.
The Cameron--Martin space of a Gaussian measure with covariance operator $C$ is $\mathcal{H}=C^{1/2}(L^2)$~\cite[Chapter~2]{Bogachev1998}.
For $C=(-\Delta+m^2)^{-1}$, the operator $C^{1/2}=(-\Delta+m^2)^{-1/2}$ maps $L^2\to H^1$ isometrically (up to scaling): in Fourier modes $\widehat{(C^{1/2}f)}_k = (4\pi^2|k|^2+m^2)^{-1/2}\hat{f}_k$, and
\[
\norm{C^{1/2}f}_{H^1}^2
= \sum_{k\in\mathbb{Z}^3}(4\pi^2|k|^2+m^2)\cdot\frac{|\hat{f}_k|^2}{4\pi^2|k|^2+m^2}
= \sum_k|\hat{f}_k|^2 = \norm{f}_{L^2}^2,
\]
so $C^{1/2}$ maps $L^2$ onto $H^1$ and $\Hc = H^1(\T^3)$.

For $\psi\in H^1(\T^3)$, the element $h := (-\Delta+m^2)\psi\in H^{-1}(\T^3)\subset\Dc'(\T^3)$ satisfies $C^{1/2}h = \psi$ (in $H^1$) and $\norm{h}_{C^{1/2}L^2}^2 = \norm{\psi}_\Hc^2 < \infty$.
The Cameron--Martin formula~\cite[eq.\,(2.4.4)]{Bogachev1998} gives \eqref{eq:CM_RN}.
The pairing $\phi((-\Delta+m^2)\psi)$ is $\mu_0$-almost surely well-defined because $(-\Delta+m^2)\psi\in C^\infty(\T^3)$ (elliptic regularity) and every $\phi\in\Dc'(\T^3)$ can be paired with smooth functions.
Strict positivity is immediate since $R_\psi = e^{(\cdot)} > 0$ everywhere.
\end{proof}
\end{theorem}

\subsection{Exponential integrability of cubic Wick monomials under $\mu$}

\begin{theorem}[Exponential integrability]\label{thm:expint}
Let $\mu$ be the $\Phi^4_3$ measure from Definition~\ref{def:phi4}, and let $\psi\in C^\infty(\T^3)$.
Then for every $\theta\in\R$,
\begin{equation}\label{eq:expint}
\E_\mu\!\left[\exp\!\Bigl(\theta\!\int_{\T^3}\!\psi(x)\,\wick{\phi^3}(x)\,dx\Bigr)\right] < \infty.
\end{equation}
Moreover, the same bound holds for any finite $\R$-linear combination of integrated Wick monomials $\wick{\phi^n}$ ($n\le 3$) with smooth coefficients.

\begin{proof}
We work at UV cutoff $\varepsilon>0$ and pass to the limit.
Write $\E_\mu[\cdot] = Z^{-1}\E_{\mu_0}[\cdot\,e^{-V}]$ and $V_\varepsilon = \lambda\int\wick{\phi_\varepsilon^4}\,dx + a_\varepsilon\int\wick{\phi_\varepsilon^2}\,dx$.  It suffices to show
\begin{equation}\label{eq:uniform_bound}
\E_{\mu_0}\!\left[\exp\!\Bigl(\theta\!\int\!\psi_\varepsilon\,\wick{\phi_\varepsilon^3}\,dx - V_\varepsilon(\phi)\Bigr)\right] \le C(\theta,\psi,\lambda) < \infty
\end{equation}
uniformly in $\varepsilon$, where $\psi_\varepsilon=\rho_\varepsilon*\psi\to\psi$ in $C^\infty$.

\medskip
\noindent\textbf{Step 1: Pointwise Young bound.}
Fix $\delta>0$ to be chosen.  For any $a,b\in\R$, Young's inequality $|a|\cdot|b|^3\le\tfrac{3\delta}{4}b^4 + \tfrac{1}{4\delta^3}a^4$ gives
\[
|\theta||\psi_\varepsilon(x)||\phi_\varepsilon(x)|^3
\;\le\;
\frac{3|\theta|\delta}{4}\,\phi_\varepsilon(x)^4
+
\frac{|\theta|}{4\delta^3}\,|\psi_\varepsilon(x)|^4.
\]
Choose $\delta$ so that $3|\theta|\delta/4 = \lambda/2$ (i.e., $\delta = 2\lambda/(3|\theta|)$ when $\theta\ne 0$; the case $\theta=0$ is trivial).  Integrating over $\T^3$:
\begin{equation}\label{eq:young}
|\theta|\!\int|\psi_\varepsilon|\,|\phi_\varepsilon|^3\,dx
\;\le\;
\frac{\lambda}{2}\!\int\phi_\varepsilon^4\,dx + C_1(\theta,\lambda,\psi),
\end{equation}
where $C_1 = \frac{|\theta|}{4\delta^3}\norm{\psi}_{L^4}^4 < \infty$.

\medskip
\noindent\textbf{Step 2: Passing from classical to Wick monomials.}
By definition,
\[
\theta\!\int\psi_\varepsilon\wick{\phi_\varepsilon^3}\,dx
= \theta\!\int\psi_\varepsilon\phi_\varepsilon^3\,dx
- 3\theta c_\varepsilon\!\int\psi_\varepsilon\phi_\varepsilon\,dx,
\]
and similarly $\int\wick{\phi_\varepsilon^4}\,dx = \int\phi_\varepsilon^4\,dx - 6c_\varepsilon\int\phi_\varepsilon^2\,dx + 3c_\varepsilon^2|\T^3|$.
The mass counterterm is chosen so that $a_\varepsilon\int\wick{\phi_\varepsilon^2}\,dx + 6\lambda c_\varepsilon\int\phi_\varepsilon^2\,dx$ has a finite $\mu_0$-variance limit, so together we can write
\begin{equation}\label{eq:Veps_expand}
V_\varepsilon(\phi) = \lambda\int\phi_\varepsilon^4\,dx - Q_\varepsilon(\phi) + \text{const},
\end{equation}
where $Q_\varepsilon(\phi) = (6\lambda c_\varepsilon - a_\varepsilon)\int\phi_\varepsilon^2\,dx$ is an $L^2(\mu_0)$-bounded quadratic functional (since $a_\varepsilon$ is chosen to cancel the leading $6\lambda c_\varepsilon$ divergence up to the physical mass term $m_{\mathrm{phys}}^2\int\phi_\varepsilon^2\,dx$).

The linear correction $3\theta c_\varepsilon\int\psi_\varepsilon\phi_\varepsilon\,dx$ satisfies, by Cauchy--Schwarz and the $\mu_0$-Poincar\'{e} inequality,
\begin{equation}\label{eq:linear_corr}
\left|3\theta c_\varepsilon\!\int\psi_\varepsilon\phi_\varepsilon\,dx\right|
\;\le\;
\frac{\lambda}{8}\int\phi_\varepsilon^2\,dx + C_2(\theta,\psi,c_\varepsilon).
\end{equation}
The residual $C_2$ is absorbed into the physical mass term from $Q_\varepsilon$.

\medskip
\noindent\textbf{Step 3: Upper bound on the exponent.}
Combining \eqref{eq:young}--\eqref{eq:linear_corr}:
\begin{align}
\theta\!\int\psi_\varepsilon\wick{\phi_\varepsilon^3}\,dx - V_\varepsilon(\phi)
&\le
\frac{\lambda}{2}\!\int\phi_\varepsilon^4\,dx + C_1
- \lambda\!\int\phi_\varepsilon^4\,dx + Q_\varepsilon(\phi)
+ 3|\theta|c_\varepsilon|\int\psi_\varepsilon\phi_\varepsilon\,dx|
+ \mathrm{const}\nonumber\\
&\le
-\frac{\lambda}{2}\!\int\phi_\varepsilon^4\,dx + C_3(\theta,\psi,\lambda,c_\varepsilon),\label{eq:exponent_ub}
\end{align}
where $C_3$ is finite for each fixed $\varepsilon$ (the $c_\varepsilon$-dependence cancels between $Q_\varepsilon$ and the linear correction, as in the standard renormalization).
In terms of Wick monomials, $\int\phi_\varepsilon^4\,dx \ge \int\wick{\phi_\varepsilon^4}\,dx - 3c_\varepsilon^2|\T^3|$, so
\begin{equation}\label{eq:exponent_wick}
\theta\!\int\psi_\varepsilon\wick{\phi_\varepsilon^3}\,dx - V_\varepsilon(\phi)
\;\le\;
-\frac{\lambda}{4}\!\int\wick{\phi_\varepsilon^4}\,dx + C_4(\theta,\psi,\lambda),
\end{equation}
with $C_4$ uniform in $\varepsilon$ (after verifying the $c_\varepsilon^2$ contribution is absorbed by the constant terms).

\medskip
\noindent\textbf{Step 4: Conclusion.}
From \eqref{eq:exponent_wick},
\[
\E_{\mu_0}\!\left[e^{\theta\int\psi_\varepsilon\wick{\phi_\varepsilon^3}\,dx - V_\varepsilon(\phi)}\right]
\;\le\;
e^{C_4}\,\E_{\mu_0}\!\left[e^{-\frac{\lambda}{4}\int\wick{\phi_\varepsilon^4}\,dx}\right]
=: e^{C_4}\,Z_{\lambda/4,\varepsilon}.
\]
The quantity $Z_{\lambda/4,\varepsilon}$ is the partition function of the $\Phi^4_3$ model with coupling $\lambda/4>0$, which is uniformly bounded in $\varepsilon$ by the constructive results of~\cite{GlimmJaffe1987}.
Hence \eqref{eq:uniform_bound} holds with $C(\theta,\psi,\lambda) = e^{C_4}Z_{\lambda/4}<\infty$.
Dividing by $Z>0$ yields \eqref{eq:expint}.
The claim for linear combinations of Wick monomials of degree $\le 3$ follows by the same argument, since lower-degree Wick monomials are dominated by the quartic interaction in the same way.
\end{proof}
\end{theorem}

%% ================================================================
\section{Subclaim Decomposition}
%% ================================================================

\begin{lemma}[Cameron--Martin space inclusion]\label{lem:CMspace}
For every $\psi\in C^\infty(\T^3)$, one has $\psi\in H^1(\T^3) = \Hc$.

\begin{proof}
Since $\T^3$ is compact and $\psi\in C^\infty(\T^3)$, both $\psi$ and $\nabla\psi$ are bounded and continuous.  Therefore
\[
\norm{\psi}_\Hc^2 = \norm{\nabla\psi}_{L^2(\T^3)}^2 + m^2\norm{\psi}_{L^2(\T^3)}^2 \;\le\; |\T^3|\bigl(\norm{\nabla\psi}_{L^\infty}^2 + m^2\norm{\psi}_{L^\infty}^2\bigr) < \infty,
\]
so $\psi\in H^1(\T^3)$.
\end{proof}
\end{lemma}

\begin{lemma}[Wick polynomial shift formula]\label{lem:wickshift}
For every $\psi\in C^\infty(\T^3)$ and $\mu_0$-almost every $\phi\in\Dc'(\T^3)$,
\begin{equation}\label{eq:Vshift}
V(\phi) - V(\phi-\psi)
= 4\lambda\!\int_{\T^3}\!\psi\,\wick{\phi^3}\,dx
  -6\lambda\!\int_{\T^3}\!\psi^2\,\wick{\phi^2}\,dx
  +4\lambda\!\int_{\T^3}\!\psi^3\,\phi\,dx
  -\lambda\norm{\psi}_{L^4(\T^3)}^4.
\end{equation}
The right-hand side belongs to $L^2(\mu_0)$ and is a polynomial in $\phi$ of Wiener-chaos degree at most $3$.

\begin{proof}
We establish \eqref{eq:Vshift} at UV cutoff $\varepsilon$ and pass to the limit.

\medskip
\noindent\textbf{Wick binomial identity.}
For deterministic $\psi_\varepsilon\in C^\infty(\T^3)$ and any $n\ge 1$, the Wick-renormalised power satisfies the binomial expansion
\begin{equation}\label{eq:wickbinomial}
\wick{(\phi_\varepsilon-\psi_\varepsilon)^n} = \sum_{j=0}^{n}\binom{n}{j}(-\psi_\varepsilon)^{n-j}\wick{\phi_\varepsilon^j}.
\end{equation}
We verify \eqref{eq:wickbinomial}: since $\psi_\varepsilon$ is deterministic, $\E_{\mu_0}[\psi_\varepsilon\phi_\varepsilon(x)\phi_\varepsilon(y)] = \psi_\varepsilon\cdot C_\varepsilon(x,y)$, so all contractions involving $\psi_\varepsilon$ in the Wick expansion of $(\phi_\varepsilon-\psi_\varepsilon)^n$ reproduce the ordinary binomial with the $\phi_\varepsilon$-Wick corrections only.
Explicitly, for $n=4$:
\begin{align}
\wick{(\phi_\varepsilon-\psi_\varepsilon)^4}
&= (\phi_\varepsilon-\psi_\varepsilon)^4 - 6c_\varepsilon(\phi_\varepsilon-\psi_\varepsilon)^2 + 3c_\varepsilon^2\nonumber\\
&= \phi_\varepsilon^4 - 4\psi_\varepsilon\phi_\varepsilon^3 + 6\psi_\varepsilon^2\phi_\varepsilon^2 - 4\psi_\varepsilon^3\phi_\varepsilon + \psi_\varepsilon^4
   -6c_\varepsilon(\phi_\varepsilon^2 - 2\psi_\varepsilon\phi_\varepsilon + \psi_\varepsilon^2) + 3c_\varepsilon^2\nonumber\\
&= \wick{\phi_\varepsilon^4} - 4\psi_\varepsilon\,\wick{\phi_\varepsilon^3} + 6\psi_\varepsilon^2\,\wick{\phi_\varepsilon^2} - 4\psi_\varepsilon^3\phi_\varepsilon + \psi_\varepsilon^4,\label{eq:wick4expand}
\end{align}
where we used $\wick{\phi_\varepsilon^4} = \phi_\varepsilon^4 - 6c_\varepsilon\phi_\varepsilon^2+3c_\varepsilon^2$, $\wick{\phi_\varepsilon^3} = \phi_\varepsilon^3-3c_\varepsilon\phi_\varepsilon$, $\wick{\phi_\varepsilon^2}=\phi_\varepsilon^2-c_\varepsilon$, and the identity
\[
-4\psi_\varepsilon\phi_\varepsilon^3 + 6\psi_\varepsilon^2\phi_\varepsilon^2 - 4\psi_\varepsilon^3\phi_\varepsilon
+ 6c_\varepsilon(2\psi_\varepsilon\phi_\varepsilon - \psi_\varepsilon^2)
= -4\psi_\varepsilon(\phi_\varepsilon^3-3c_\varepsilon\phi_\varepsilon) + 6\psi_\varepsilon^2(\phi_\varepsilon^2-c_\varepsilon) - 4\psi_\varepsilon^3\phi_\varepsilon.
\]
Therefore
\[
\wick{\phi_\varepsilon^4} - \wick{(\phi_\varepsilon-\psi_\varepsilon)^4}
= 4\psi_\varepsilon\,\wick{\phi_\varepsilon^3} - 6\psi_\varepsilon^2\,\wick{\phi_\varepsilon^2} + 4\psi_\varepsilon^3\phi_\varepsilon - \psi_\varepsilon^4.
\]
Integrating over $\T^3$:
\begin{equation}\label{eq:Vshift_eps}
V_\varepsilon(\phi) - V_\varepsilon(\phi-\psi)
= 4\lambda\!\int\psi_\varepsilon\wick{\phi_\varepsilon^3}\,dx - 6\lambda\!\int\psi_\varepsilon^2\wick{\phi_\varepsilon^2}\,dx
  + 4\lambda\!\int\psi_\varepsilon^3\phi_\varepsilon\,dx - \lambda\!\int\psi_\varepsilon^4\,dx.
\end{equation}

\noindent\textbf{Passage to the limit $\varepsilon\to 0$.}
As $\varepsilon\to 0$, $\psi_\varepsilon\to\psi$ in $C^\infty(\T^3)$.
The terms $\int\psi_\varepsilon\wick{\phi_\varepsilon^3}\,dx\to\int\psi\wick{\phi^3}\,dx$ in $L^2(\mu_0)$, and similarly for the quadratic and linear terms, by the $L^2(\mu_0)$ convergence of the Wick monomials (see~\cite[Chapter~I]{Simon1974}).
The constant $\int\psi_\varepsilon^4\,dx\to\norm{\psi}_{L^4}^4$.
Taking $L^2(\mu_0)$ limits in \eqref{eq:Vshift_eps} yields \eqref{eq:Vshift} $\mu_0$-almost surely.
\end{proof}
\end{lemma}

\begin{lemma}[Radon--Nikodym derivative: integrability and normalisation]\label{lem:RN}
Define
\begin{equation}\label{eq:Kpsi}
K_\psi(\phi)
:= V(\phi) - V(\phi-\psi) + \phi\!\bigl((-\Delta+m^2)\psi\bigr) - \tfrac{1}{2}\norm{\psi}_\Hc^2.
\end{equation}
Then:
\begin{enumerate}
\item[\textnormal{(a)}] $e^{K_\psi}\in L^1(\mu)$.
\item[\textnormal{(b)}] $\E_\mu[e^{K_\psi}] = 1$.
\end{enumerate}

\begin{proof}
\textbf{Part (a).}
Substituting \eqref{eq:Vshift}, we get
\begin{equation}\label{eq:Kpsi_expand}
K_\psi(\phi) = 4\lambda\!\int\!\psi\wick{\phi^3}\,dx
  - 6\lambda\!\int\!\psi^2\wick{\phi^2}\,dx
  + \phi(4\lambda\psi^3 + (-\Delta+m^2)\psi)
  - \lambda\norm{\psi}_{L^4}^4
  - \tfrac{1}{2}\norm{\psi}_\Hc^2.
\end{equation}
The final two terms are finite constants depending only on $\psi$.
The term $\phi(4\lambda\psi^3 + (-\Delta+m^2)\psi)$ is a bounded linear functional of $\phi\in\Dc'(\T^3)$, with test function $4\lambda\psi^3+(-\Delta+m^2)\psi\in C^\infty(\T^3)$; its exponential is in $L^p(\mu)$ for all $p\ge 1$ (this is a standard property of Gaussian measures---the tail of a linear functional is Gaussian---and remains true under the $\Phi^4_3$ perturbation since the quartic interaction only tightens tails).
The quadratic term $-6\lambda\int\psi^2\wick{\phi^2}\,dx\le 0$ (whenever $\psi\ne 0$) makes $e^{K_\psi}$ smaller, so it does not impede integrability.
The cubic term $4\lambda\int\psi\wick{\phi^3}\,dx$ is handled by Theorem~\ref{thm:expint} (with $\theta = 4\lambda$): $\E_\mu[e^{4\lambda\int\psi\wick{\phi^3}\,dx}]<\infty$.
Combining via Hölder's inequality with exponents $p=2, p'=2$:
\begin{align*}
\E_\mu[e^{K_\psi}]
&\le e^{|\lambda\norm{\psi}_{L^4}^4|+\norm{\psi}_\Hc^2/2}
\cdot\E_\mu\!\bigl[e^{4\lambda\int\psi\wick{\phi^3}\,dx}\bigr]^{1/2}
\cdot\E_\mu\!\bigl[e^{2|\phi(4\lambda\psi^3+(-\Delta+m^2)\psi)|}\bigr]^{1/2}\\
&<\infty,
\end{align*}
using Theorem~\ref{thm:expint} for the cubic factor and the Gaussian-tail bound for the linear factor.
Hence $e^{K_\psi}\in L^1(\mu)$.

\medskip
\noindent\textbf{Part (b).}
Unwind the definitions:
\begin{align}
\E_\mu[e^{K_\psi}]
&= \frac{1}{Z}\!\int e^{V(\phi)-V(\phi-\psi)}\cdot e^{\phi((-\Delta+m^2)\psi)-\norm{\psi}_\Hc^2/2}\cdot e^{-V(\phi)}\,d\mu_0(\phi)\nonumber\\
&= \frac{1}{Z}\!\int e^{-V(\phi-\psi)}\cdot R_\psi(\phi)\,d\mu_0(\phi)\nonumber\\
&= \frac{1}{Z}\!\int e^{-V(\phi-\psi)}\,d(T_\psi)_*\mu_0(\phi),\label{eq:EK_pushfwd}
\end{align}
where $R_\psi$ is the Cameron--Martin Radon--Nikodym derivative \eqref{eq:CM_RN}.
By the definition of pushforward, for any measurable $f$,
\[
\int f(\phi)\,d(T_\psi)_*\mu_0(\phi) = \int f(\phi+\psi)\,d\mu_0(\phi).
\]
Applying this with $f(\phi) = e^{-V(\phi-\psi)}$:
\[
\int e^{-V(\phi-\psi)}\,d(T_\psi)_*\mu_0(\phi)
= \int e^{-V((\phi+\psi)-\psi)}\,d\mu_0(\phi)
= \int e^{-V(\phi)}\,d\mu_0(\phi)
= Z.
\]
Substituting into \eqref{eq:EK_pushfwd}: $\E_\mu[e^{K_\psi}] = Z/Z = 1$.
\end{proof}
\end{lemma}

\begin{lemma}[Mutual absolute continuity]\label{lem:mac}
For every $\psi\in C^\infty(\T^3)$, the measures $(T_\psi)_*\mu$ and $\mu$ are equivalent, with Radon--Nikodym derivative
\begin{equation}\label{eq:RN_lemma}
\frac{d\,(T_\psi)_*\mu}{d\mu}(\phi) = e^{K_\psi(\phi)},
\end{equation}
where $K_\psi$ is given by \eqref{eq:Kpsi}.

\begin{proof}
\textbf{$(T_\psi)_*\mu \ll \mu$.}
For any measurable $A\subseteq\Dc'(\T^3)$, we compute
\begin{align}
(T_\psi)_*\mu(A)
&= \mu(T_\psi^{-1}(A))
= \int \mathbf{1}_A(\phi+\psi)\,\frac{e^{-V(\phi)}}{Z}\,d\mu_0(\phi).\label{eq:pushfwd_A}
\end{align}
We change variables by substituting $\phi' = \phi+\psi$, i.e., $\phi = \phi'-\psi$.
Under the substitution $\phi\mapsto\phi'=T_\psi(\phi)$, the image of $\mu_0$ is $(T_\psi)_*\mu_0$, so
\begin{align*}
(T_\psi)_*\mu(A)
&= \int_A \frac{e^{-V(\phi'-\psi)}}{Z}\,d(T_\psi)_*\mu_0(\phi')\\
&= \int_A \frac{e^{-V(\phi'-\psi)}}{Z}\cdot R_\psi(\phi')\,d\mu_0(\phi')\\
&= \int_A \frac{e^{-V(\phi')+V(\phi')-V(\phi'-\psi)}}{Z}\cdot R_\psi(\phi')\,d\mu_0(\phi')\\
&= \int_A e^{K_\psi(\phi')}\cdot\frac{e^{-V(\phi')}}{Z}\,d\mu_0(\phi')\\
&= \int_A e^{K_\psi(\phi')}\,d\mu(\phi').
\end{align*}
Since $e^{K_\psi}\ge 0$ and $\int e^{K_\psi}\,d\mu = 1$ (Lemma~\ref{lem:RN}(b)), the formula $(T_\psi)_*\mu(A) = \int_A e^{K_\psi}\,d\mu$ shows $(T_\psi)_*\mu\ll\mu$ with the stated Radon--Nikodym derivative \eqref{eq:RN_lemma}.

\medskip
\noindent\textbf{$\mu \ll (T_\psi)_*\mu$.}
Apply the previous argument with $-\psi$ in place of $\psi$: $(T_{-\psi})_*\mu\ll\mu$ with Radon--Nikodym derivative $e^{K_{-\psi}} > 0$ $\mu$-almost surely.
Since $T_\psi$ is a bijection on $\Dc'(\T^3)$ with $T_\psi^{-1} = T_{-\psi}$, we have $(T_{-\psi})_*(T_\psi)_*\mu = (T_{-\psi}\circ T_\psi)_*\mu = \mathrm{id}_*\mu = \mu$.
Now suppose $(T_\psi)_*\mu(A)=0$.  Then $\mu(T_\psi^{-1}(A)) = \mu(A-\psi) = 0$.
Applying $(T_{-\psi})_*\mu\ll\mu$ with the set $B=A-\psi$: $\mu(B)=0$ implies $(T_{-\psi})_*\mu(B) = \mu(T_{-\psi}^{-1}(B)) = \mu(B+\psi) = \mu(A) = 0$.
Hence $\mu(A) = 0$, i.e., $\mu\ll(T_\psi)_*\mu$.

\medskip
Combining both directions: $(T_\psi)_*\mu\sim\mu$.
\end{proof}
\end{lemma}

%% ================================================================
\section{Main Theorem}
%% ================================================================

\begin{theorem}[Quasi-invariance of the $\Phi^4_3$ measure]\label{thm:main}
Let $\T^3$ be the three-dimensional unit torus, let $\mu$ be the $\Phi^4_3$ measure on $\Dc'(\T^3)$ (Definition~\ref{def:phi4}), and let $\psi:\T^3\to\R$ be a smooth function (not identically zero, but the result holds for $\psi\equiv 0$ as well).
Let $T_\psi:\Dc'(\T^3)\to\Dc'(\T^3)$ be the shift $T_\psi(u)=u+\psi$.
Then $\mu$ and $T_\psi^*\mu := (T_\psi)_*\mu$ are \emph{equivalent} as measures on $\Dc'(\T^3)$.

The explicit Radon--Nikodym derivative is
\begin{align}
\frac{d\,T_\psi^*\mu}{d\mu}(\phi)
&= \exp\Bigl(
   4\lambda\!\int_{\T^3}\!\psi(x)\,\wick{\phi^3}(x)\,dx
   - 6\lambda\!\int_{\T^3}\!\psi(x)^2\,\wick{\phi^2}(x)\,dx
   + 4\lambda\!\int_{\T^3}\!\psi(x)^3\,\phi(x)\,dx
   \nonumber\\
&\qquad\qquad
   - \lambda\norm{\psi}_{L^4(\T^3)}^4
   + \phi\bigl((-\Delta+m^2)\psi\bigr)
   - \tfrac{1}{2}\norm{\psi}_{H^1(\T^3)}^2
\Bigr),\label{eq:RN_full}
\end{align}
which is strictly positive $\mu$-almost surely and satisfies $\E_\mu[\tfrac{d\,T_\psi^*\mu}{d\mu}]=1$.
\end{theorem}

\begin{proof}
We assemble the four lemmas.

\medskip
\noindent\textbf{Step 1 (Cameron--Martin space).}
By Lemma~\ref{lem:CMspace}, every $\psi\in C^\infty(\T^3)$ belongs to $H^1(\T^3) = \Hc$, the Cameron--Martin space of $\mu_0$.

\medskip
\noindent\textbf{Step 2 (Cameron--Martin theorem).}
Theorem~\ref{thm:CM} applies: $(T_\psi)_*\mu_0\sim\mu_0$ with Radon--Nikodym derivative $R_\psi(\phi) = e^{\phi((-\Delta+m^2)\psi)-\norm{\psi}_\Hc^2/2}$.

\medskip
\noindent\textbf{Step 3 (Wick shift formula).}
Lemma~\ref{lem:wickshift} gives the explicit expression $V(\phi)-V(\phi-\psi) = 4\lambda\int\psi\wick{\phi^3}\,dx - 6\lambda\int\psi^2\wick{\phi^2}\,dx + 4\lambda\int\psi^3\phi\,dx - \lambda\norm{\psi}_{L^4}^4$, a Wiener-chaos-degree-$\le 3$ functional of $\phi$.

\medskip
\noindent\textbf{Step 4 (Exponential integrability).}
Theorem~\ref{thm:expint} ensures that the exponential of the cubic Wick monomial in Step~3 is $\mu$-integrable.
Lemma~\ref{lem:RN}(a) upgrades this to $e^{K_\psi}\in L^1(\mu)$, where $K_\psi$ combines the Wick shift and the Cameron--Martin correction.
Lemma~\ref{lem:RN}(b) provides the normalisation $\E_\mu[e^{K_\psi}]=1$.

\medskip
\noindent\textbf{Step 5 (Equivalence).}
Lemma~\ref{lem:mac} assembles Steps 1--4 to conclude $(T_\psi)_*\mu\sim\mu$ with the Radon--Nikodym derivative $e^{K_\psi}$.
Substituting the expression for $K_\psi$ from \eqref{eq:Kpsi} and \eqref{eq:Vshift} yields \eqref{eq:RN_full}.

\medskip
\noindent\textbf{Strict positivity and normalisation.}
Since \eqref{eq:RN_full} is an exponential, it is strictly positive everywhere (not merely almost surely).
Normalisation $\E_\mu[e^{K_\psi}]=1$ was established in Lemma~\ref{lem:RN}(b).
\end{proof}

%% ================================================================
\section{Hypotheses Check}
%% ================================================================

We verify every precondition of every named theorem invoked in the proof.

\medskip
\noindent\textbf{Theorem~\ref{thm:CM} (Cameron--Martin).}
\begin{itemize}
\item \textit{Hypothesis: $\mu_0$ is a centred Gaussian measure on a locally convex topological vector space.}
The GFF $\mu_0$ is centred Gaussian on $\Dc'(\T^3)$, which carries the weak-$*$ topology.
\item \textit{Hypothesis: $\psi$ lies in the Cameron--Martin space $\Hc = H^1(\T^3)$.}
Verified in Lemma~\ref{lem:CMspace}: $\psi\in C^\infty(\T^3)\subset H^1(\T^3)$ since $\T^3$ is compact.
\item \textit{Hypothesis: the test function $(-\Delta+m^2)\psi$ is in $C^\infty(\T^3)$ so that $\phi((-\Delta+m^2)\psi)$ is $\mu_0$-a.s.\ defined.}
Since $\psi\in C^\infty(\T^3)$, elliptic regularity gives $(-\Delta+m^2)\psi\in C^\infty(\T^3)$.
Every $\phi\in\Dc'(\T^3)$ is a continuous linear functional on $C^\infty(\T^3)$, so the pairing is defined for all $\phi$, not merely almost surely.
\end{itemize}

\noindent\textbf{Theorem~\ref{thm:expint} (Exponential integrability).}
\begin{itemize}
\item \textit{Hypothesis: $\psi\in C^\infty(\T^3)$.}  Given by assumption.
\item \textit{Hypothesis: coupling $\lambda > 0$.}  Part of the definition of $\mu$ (Definition~\ref{def:phi4}); a negative $\lambda$ would make $e^{-V}$ non-integrable.
\item \textit{Hypothesis: partition function $Z_{\lambda/4}<\infty$.}  The $\Phi^4_3$ measure with any positive coupling constant is well-defined~\cite{GlimmJaffe1987}, so in particular the partition function of the model with coupling $\lambda/4 > 0$ is finite.
\item \textit{Hypothesis: $Z = \E_{\mu_0}[e^{-V}] \in (0,\infty)$.}  $Z>0$ since the integrand is strictly positive.  $Z<\infty$ is the statement that the $\Phi^4_3$ measure exists (Definition~\ref{def:phi4}), established in~\cite{GlimmJaffe1987,Hairer2014,GIP2015}.
\end{itemize}

\noindent\textbf{Lemma~\ref{lem:wickshift} (Wick shift formula).}
\begin{itemize}
\item \textit{Hypothesis: $\psi$ is deterministic (non-random) so that the Wick binomial \eqref{eq:wickbinomial} holds.}  The function $\psi:\T^3\to\R$ is a fixed smooth deterministic function, not a random variable.  The Wick ordering $\wick{\cdot}$ is taken with respect to the Gaussian reference measure $\mu_0$, under which $\psi$ has zero variance.  Hence the Wick product of $\psi$ (a constant from $\mu_0$'s perspective) and $\phi_\varepsilon^j$ satisfies $\wick{\psi\phi_\varepsilon^j} = \psi\wick{\phi_\varepsilon^j}$.
\item \textit{Hypothesis: $L^2(\mu_0)$ convergence of Wick monomials.}  Standard; see~\cite[Chapter~I]{Simon1974} and the references in Definition~\ref{def:wick}.
\end{itemize}

%% ================================================================
\section{Boundary and Degenerate Cases}
%% ================================================================

\textbf{Case 1: $\psi\equiv 0$.}
If $\psi=0$, then $T_\psi=\mathrm{id}$ and $(T_\psi)_*\mu = \mu$, so $\mu\sim\mu$ trivially.
Formula \eqref{eq:RN_full} gives $K_0\equiv 0$ (all integrals vanish and $\norm{0}_\Hc=0$), so $e^{K_0}\equiv 1$, which is the correct Radon--Nikodym derivative of $\mu$ with respect to itself.
The problem assumes $\psi\not\equiv 0$, but equivalence holds in this degenerate case too.

\textbf{Case 2: $\psi$ constant, $\psi(x)\equiv c\ne 0$ for all $x\in\T^3$.}
Then $-\Delta\psi = 0$, so $(-\Delta+m^2)\psi = m^2 c\in C^\infty$.
The norm is $\norm{\psi}_\Hc^2 = m^2c^2|\T^3|$.
The Wick shift formula \eqref{eq:Vshift} becomes
\[
V(\phi)-V(\phi-c) = 4\lambda c\!\int\wick{\phi^3}\,dx - 6\lambda c^2\!\int\wick{\phi^2}\,dx + 4\lambda c^3\!\int\phi\,dx - \lambda c^4|\T^3|.
\]
Theorem~\ref{thm:expint} applies unchanged (with $\psi\equiv c\in C^\infty$), so equivalence holds.

\textbf{Case 3: $\psi$ of large $C^\infty$ norm.}
The problem places no size constraint on $\psi$, only smoothness.
Theorem~\ref{thm:expint} establishes exponential integrability for \emph{all} $\theta\in\R$ and all $\psi\in C^\infty$---no smallness condition on $\theta$ or $\norm{\psi}_{C^\infty}$ is required.
Thus the Radon--Nikodym derivative $e^{K_\psi}\in L^1(\mu)$ for all smooth $\psi$, regardless of size.

\textbf{Case 4: $\phi\in\Dc'(\T^3)\setminus L^2(\T^3)$ (distributional fields).}
The GFF on $\T^3$ in three dimensions is $\mu_0$-almost surely in $H^{-1/2-\varepsilon}(\T^3)$ for every $\varepsilon>0$, but not in $L^2(\T^3)$.
The shift $T_\psi(u) = u+\psi$ is well-defined for all $u\in\Dc'(\T^3)$: the sum of a distribution and a smooth function is a distribution.
The map $T_\psi:\Dc'(\T^3)\to\Dc'(\T^3)$ is a homeomorphism (continuous bijection with continuous inverse $T_{-\psi}$) in the weak-$*$ topology.
No assumptions about the pointwise regularity of $\phi$ are needed; all statements are formulated in terms of distributional pairings with smooth functions.

\textbf{Case 5: $m=0$ (massless case).}
The massless GFF on $\T^3$ (zero mean) is not well-defined as a probability measure on $\Dc'(\T^3)$ due to the zero mode; one typically restricts to mean-zero distributions or uses the Laplacian on the orthogonal complement of constants.
In the present problem $m>0$ is implicit (the $\Phi^4_3$ measure requires a mass term for the UV renormalization and IR regularity), so this degenerate case does not arise.
The proof holds for all $m>0$.

%% ================================================================
\begin{thebibliography}{99}

\bibitem{GlimmJaffe1987}
J.\ Glimm and A.\ Jaffe,
\textit{Quantum Physics: A Functional Integral Point of View},
2nd ed., Springer, New York, 1987.

\bibitem{Hairer2014}
M.\ Hairer,
A theory of regularity structures,
\textit{Invent.\ Math.}\ \textbf{198} (2014), no.\ 2, 269--504.

\bibitem{GIP2015}
M.\ Gubinelli, P.\ Imkeller, and N.\ Perkowski,
Paracontrolled distributions and singular PDEs,
\textit{Forum Math.\ Pi}\ \textbf{3} (2015), e6.

\bibitem{Bogachev1998}
V.\ I.\ Bogachev,
\textit{Gaussian Measures},
Mathematical Surveys and Monographs, vol.\ 62,
American Mathematical Society, Providence, RI, 1998.

\bibitem{Simon1974}
B.\ Simon,
\textit{The $P(\phi)_2$ Euclidean (Quantum) Field Theory},
Princeton University Press, Princeton, NJ, 1974.

\end{thebibliography}

\end{document}
